% ============================================================
%  SECTION 3.3 — CORRECTED AND IMPROVED VERSION
%  Representational Ethical Space & Ethical Effective Hyper-Volume
%  Author: M. Azzano — revised with structural corrections
% ============================================================

\subsection{The Representational Ethical Space and the Ethical Effective Hyper-Volume}
\label{sec:rep-ethical-space}

Let $\Omega$ denote the total phase space of the world-model.  
An agent~$a$ at time~$t$ maintains an \textbf{internal representational subspace}
\begin{equation}
    \Omega_a(t) \subseteq \Omega,
\end{equation}
called the \textbf{representational ethical space} of the agent.

The space $\Omega_a(t)$ is characterised by:
\begin{itemize}
    \item \emph{Dimensionality}\; $d_a(t)=\dim\bigl(\Omega_a(t)\bigr)$\,;
    \item \emph{Riemannian volume}\; $V_a(t)=\int_{\Omega_a(t)}\!\sqrt{\det g_{ij}(q)}\;d^{\,n}q$\,,
          where $g_{ij}$ is the Fisher--Rao metric induced by the agent's internal probabilistic model;
    \item \emph{Temporal depth}\; $\tau_a(t)$, the horizon over which the agent can propagate
          consequential predictions (defined precisely below).
\end{itemize}

\subsubsection{Characteristic prediction time}
\label{sec:tau-char}

For each point $q\in\Omega_a$, define the \textbf{regularised mutual information}
\begin{equation}
    I_\varepsilon\bigl(\mathbf{x}(t);\,\mathbf{x}(t+\tau)\,\big|\,q\bigr)
    \;=\;
    I\!\left(\mathbf{x}^{(\varepsilon)}(t);\,\mathbf{x}^{(\varepsilon)}(t+\tau)\,\Big|\,q\right),
\end{equation}
where $\mathbf{x}^{(\varepsilon)}$ is the state coarse-grained at resolution~$\varepsilon$
(e.g.\ $\varepsilon$-binning or Gaussian convolution with bandwidth~$\varepsilon$).
This regularisation ensures that $I_\varepsilon$ is finite and differentiable at $\tau=0$.

The \textbf{characteristic prediction time} at~$q$ is then
\begin{equation}
\label{eq:tau-char}
    \tau_{\mathrm{char}}(q)
    \;=\;
    -\left[\frac{d}{d\tau}\log I_\varepsilon
       \bigl(\mathbf{x}(t);\,\mathbf{x}(t+\tau)\,\big|\,q\bigr)
    \right]^{-1}_{\!\tau=0}\,.
\end{equation}
Physical requirement: $\tau_{\mathrm{char}}(q)>0$, which holds whenever the
agent's internal model has strictly decaying predictability
(guaranteed for ergodic Markov models; cf.\ the data-processing inequality).

\begin{remark}[Interpretation]
$\tau_{\mathrm{char}}(q)$ is the $e$-folding time of predictive mutual information
at configuration~$q$.  Points where the agent can forecast far into the
future ($\tau_{\mathrm{char}}$ large) contribute more to ethical capacity;
points of rapid decoherence contribute less.
\end{remark}

\subsubsection{Ethical effective hyper-volume}
\label{sec:phi}

\begin{definition}[Ethical effective hyper-volume]
\begin{equation}
\label{eq:phi}
\boxed{\;
    \Phi_a(t)
    \;=\;
    \int_{\Omega_a(t)}\!
        \sqrt{\det g_{ij}(q)}\;\tau_{\mathrm{char}}(q)\;d^{\,n}q\,.
\;}
\end{equation}
\end{definition}

\begin{remark}[Derivation from the temporal weight kernel]
The definition~\eqref{eq:phi} arises from integrating the exponential
temporal kernel $\rho(\tau;q)=\exp\!\bigl(-\tau/\tau_{\mathrm{char}}(q)\bigr)$
over the full temporal axis:
\[
    \Phi_a(t)
    \;=\;
    \int_{\Omega_a(t)}\!\sqrt{\det g_{ij}(q)}
    \left[\int_0^{\infty}\!\exp\!\Bigl(-\frac{\tau}{\tau_{\mathrm{char}}(q)}\Bigr)d\tau\right]
    d^{\,n}q
    \;=\;
    \int_{\Omega_a(t)}\!\sqrt{\det g_{ij}(q)}\;\tau_{\mathrm{char}}(q)\;d^{\,n}q\,.
\]
The temporal dimension is thus \emph{analytically integrated out},
leaving a Riemannian volume weighted by local predictive horizon.
\end{remark}


% ============================================================
\subsubsection{Coupling to the Ethical Lagrangian}
\label{sec:coupling}

The growth rate of integrated consciousness information satisfies the
\textbf{operative capacity constraint}
\begin{equation}
\label{eq:capacity-chain}
    \dot{I}_c(t)
    \;\leq\;
    f\!\bigl(\Phi_a^{(\mathrm{used})}(t)\bigr)
    \;\leq\;
    f\!\bigl(\Phi_a(t)\bigr),
\end{equation}
where $\Phi_a^{(\mathrm{used})}(t)\leq\Phi_a(t)$ is the ethical hyper-volume
that the agent \emph{actually deploys} in the decision at time~$t$.
The first inequality encodes the fact that an agent can integrate
information only through the representational volume it actively
employs; the second follows from monotonicity of~$f$.

\paragraph{Why the constraint must involve $\Phi_a^{(\mathrm{used})}$, not $\Phi_a$.}
The total hyper-volume $\Phi_a$ measures \emph{capacity};
the operative bound on $\dot{I}_c$ must reflect \emph{actual deployment}.
A library one owns but never opens does not improve one's decisions.
This distinction is the formal root of the concept of ethical negligence
(Section~\ref{sec:responsibility}).

\subsubsection{Emergent form of the bounding function $f(\Phi)$}
\label{sec:f-derivation}

Let $\hat{\mathcal{F}}$ be the Fisher information operator on the
tangent bundle $T\Omega_a$, with eigenvalues $\{\lambda_k\}_{k=1}^{d_a}$
ordered decreasingly.  Define the \textbf{total Fisher capacity}
\begin{equation}
    \Lambda_a(t)=\operatorname{Tr}\hat{\mathcal{F}}
    =\sum_{k=1}^{d_a(t)}\lambda_k(t),
\end{equation}
and the \textbf{dominant eigenvalue} $\lambda_{\max}(t)=\max_k\lambda_k(t)$.

\begin{proposition}[Channel-capacity bound]
\label{prop:f-form}
The bounding function has the saturating form
\begin{equation}
\label{eq:f-form}
\boxed{\;
    f(\Phi)=\alpha\,\log\!\bigl(1+\beta\,\Phi\bigr),
    \qquad
    \alpha=\frac{\Lambda_a}{2\pi},
    \qquad
    \beta=\frac{\lambda_{\max}}{\Lambda_a}\,.
\;}
\end{equation}
Both constants are determined solely by the Fisher spectrum of the
agent's internal model; no external hyperparameters are introduced.
\end{proposition}

\begin{proof}
The agent's internal model defines a communication channel between the
world state and the internal representation.  Each eigenmode~$k$ of
$\hat{\mathcal{F}}$ acts as an independent Gaussian sub-channel whose
signal-to-noise ratio is proportional to $\lambda_k\,\Phi_k$, where
$\Phi_k$ is the hyper-volume projected onto the $k$-th eigenmode.

By the continuous-time Gaussian channel capacity theorem
\cite{Shannon1948,CoverThomas2006}, the capacity per unit time of a
single mode with spectral density $\lambda_k$ and effective power
$\Phi_k$ is
\[
    C_k = \frac{1}{2}\log(1+\lambda_k\,\Phi_k)\,.
\]
The aggregate capacity, summing over modes and applying the water-filling
bound with uniform volume allocation $\Phi_k\approx\Phi/d_a$, yields
\[
    C_{\mathrm{tot}}
    \;\leq\;
    \sum_{k=1}^{d_a}\frac{1}{2}\log\!\Bigl(1+\frac{\lambda_k}{d_a}\,\Phi\Bigr)
    \;\leq\;
    \frac{d_a}{2}\,\log\!\Bigl(1+\frac{\lambda_{\max}}{d_a}\,\Phi\Bigr),
\]
where the last inequality replaces each $\lambda_k$ by $\lambda_{\max}$.
Rewriting $d_a/2 = \Lambda_a/(2\bar\lambda)$ where
$\bar\lambda=\Lambda_a/d_a$ is the mean eigenvalue, and absorbing
the conversion to a continuous-time rate (which introduces the
$2\pi$ normalisation from the spectral density of the Nyquist
bandwidth), one obtains
\[
    \dot{I}_c
    \;\leq\;
    \frac{\Lambda_a}{2\pi}\,\log\!\Bigl(1+\frac{\lambda_{\max}}{\Lambda_a}\,\Phi\Bigr)
    \;=\;
    \alpha\,\log(1+\beta\,\Phi).
    \qedhere
\]
\end{proof}

\begin{remark}[Properties of $f$]
$f$ is strictly increasing, strictly concave, and $f(0)=0$.
These properties are essential for Theorem~\ref{thm:lyapunov}.
The ratio $\beta=\lambda_{\max}/\Lambda_a\in(0,1]$ measures
spectral concentration: $\beta=1$ for a single dominant mode,
$\beta\to 0$ for a flat spectrum.
\end{remark}

Substituting into the EECF Lagrangian:
\begin{equation}
\label{eq:Leth-constrained}
    \mathcal{L}_{\mathrm{eth}}
    = \dot{I}_c - \dot{S}_{\mathrm{ent}}
    \quad\text{subject to}\quad
    \dot{I}_c
    \leq \frac{\Lambda_a}{2\pi}
         \log\!\Bigl(1+\frac{\lambda_{\max}}{\Lambda_a}\,\Phi_a^{(\mathrm{used})}\Bigr).
\end{equation}


% ============================================================
\subsubsection{Ethical Responsibility as Unused Representational Capacity}
\label{sec:responsibility}

\begin{definition}[Ethical negligence]
\begin{equation}
\label{eq:Ra}
    R_a(t)
    \;=\;
    \Phi_a(t)-\Phi_a^{(\mathrm{used})}(t)
    \;\geq 0.
\end{equation}
$R_a$ measures the discrepancy between the map the agent \emph{could}
deploy and the map it \emph{chose} to deploy.  An agent with $R_a>0$
is formally \textbf{ethically negligent}: it possesses representational
capacity that, if deployed, would raise the ceiling on $\dot{I}_c$
and hence on $\mathcal{L}_{\mathrm{eth}}$.
\end{definition}


% ============================================================
\subsubsection{Ethical Convergence: $R_a$ as a Lyapunov-like Functional}
\label{sec:lyapunov}

\begin{assumption}[Gradient-following optimiser]
\label{ass:gradient}
The agent updates its deployed volume $\Phi_a^{(\mathrm{used})}$ by
(approximate) gradient ascent on $\mathcal{L}_{\mathrm{eth}}$ with
respect to its internal allocation, with an effective learning rate
$\gamma>0$.
\end{assumption}

\begin{assumption}[Non-degenerative learning]
\label{ass:nondeg}
$\dot{\Phi}_a(t)\geq 0$: the agent's total representational capacity
does not contract.
\end{assumption}

\begin{remark}
Assumption~\ref{ass:nondeg} can be violated (neurodegeneration,
catastrophic forgetting, externally imposed information destruction).
The case $\dot{\Phi}_a<0$ is treated separately in
Corollary~\ref{cor:censorship}.
\end{remark}

\begin{theorem}[Lyapunov stability of ethical convergence]
\label{thm:lyapunov}
Under Assumptions~\ref{ass:gradient}--\ref{ass:nondeg}, $R_a(t)$ is
a Lyapunov-like functional:
\begin{equation}
\label{eq:lyapunov}
    \frac{d}{dt}R_a(t)\;\leq\; 0
\end{equation}
along system trajectories, with equality if and only if
$\Phi_a^{(\mathrm{used})}(t)=\Phi_a(t)$.
Moreover, $R_a(t)\to 0$ asymptotically.
\end{theorem}

\begin{proof}
\textbf{Step 1 (Gradient pressure on deployed volume).}
Since the operative constraint is
$\dot{I}_c\leq f\bigl(\Phi_a^{(\mathrm{used})}\bigr)$,
the Lagrangian $\mathcal{L}_{\mathrm{eth}}$ is maximised when
$\Phi_a^{(\mathrm{used})}$ is as large as possible (because $f$ is
strictly increasing).  Under Assumption~\ref{ass:gradient}, the
deployed volume evolves as
\begin{equation}
\label{eq:phi-used-dynamics}
    \dot{\Phi}_a^{(\mathrm{used})}
    \;=\;
    \dot{\Phi}_a
    \;+\;
    \gamma\,f'\!\bigl(\Phi_a^{(\mathrm{used})}\bigr)\,
    \bigl(\Phi_a - \Phi_a^{(\mathrm{used})}\bigr),
\end{equation}
where the first term reflects passive growth (the used volume inherits
the expansion of the total space) and the second term reflects
active optimisation pressure pushing $\Phi_a^{(\mathrm{used})}$
toward $\Phi_a$.  The factor $f'(\Phi_a^{(\mathrm{used})})>0$ arises
as the marginal return on deploying additional representational volume.

\textbf{Step 2 (Lyapunov inequality).}
Differentiate $R_a=\Phi_a-\Phi_a^{(\mathrm{used})}$:
\begin{equation}
    \dot{R}_a
    = \dot{\Phi}_a - \dot{\Phi}_a^{(\mathrm{used})}
    = -\gamma\,f'\!\bigl(\Phi_a^{(\mathrm{used})}\bigr)\,R_a\,.
\end{equation}
Define
\begin{equation}
    \kappa(t)
    \;=\;
    \gamma\,f'\!\bigl(\Phi_a^{(\mathrm{used})}(t)\bigr)
    \;=\;
    \gamma\,\frac{\alpha\,\beta}{1+\beta\,\Phi_a^{(\mathrm{used})}(t)}
    \;>\; 0\,.
\end{equation}
Then
\begin{equation}
\label{eq:Rdot}
    \dot{R}_a(t) = -\kappa(t)\,R_a(t) \;\leq\; 0,
\end{equation}
with equality if and only if $R_a=0$.

\textbf{Step 3 (Asymptotic convergence).}
Integrating~\eqref{eq:Rdot}:
\[
    R_a(t)
    = R_a(0)\,\exp\!\Bigl(-\int_0^t\!\kappa(s)\,ds\Bigr).
\]
Since $\kappa(t)$ is bounded below by
$\kappa_{\min}=\gamma\alpha\beta/(1+\beta\,\Phi_a^{(\max)})>0$
whenever $\Phi_a$ remains bounded, we obtain exponential decay:
\[
    R_a(t)\leq R_a(0)\,e^{-\kappa_{\min}\,t}\;\longrightarrow\; 0
    \quad\text{as }t\to\infty.
    \qedhere
\]
\end{proof}

\begin{remark}[Compactness]
The exponential-decay argument replaces the LaSalle invariance
principle used in the preliminary version, avoiding the need for
compact sublevel sets in possibly infinite-dimensional spaces.
\end{remark}


% ============================================================
\subsubsection{Corollaries: Censorship, Negligence, and Architectural Implications}
\label{sec:corollaries}

\begin{corollary}[Anti-ethical nature of capacity reduction]
\label{cor:censorship}
Let an external perturbation reduce $\Phi_a$ by $\Delta\Phi>0$
(censorship, coercion, information destruction).  Then:
\begin{enumerate}
    \item[(i)]  The capacity ceiling drops:
        $f(\Phi_a)\to f(\Phi_a-\Delta\Phi)<f(\Phi_a)$.
    \item[(ii)] The maximum achievable $\mathcal{L}_{\mathrm{eth}}$ decreases.
    \item[(iii)] If the reduction affects $\Phi_a$ but not
        $\Phi_a^{(\mathrm{used})}$ (the agent was already under-deploying),
        then $R_a$ decreases, but the \emph{ceiling on future ethical
        growth} is permanently lowered.
\end{enumerate}
Therefore, external capacity reduction is \textbf{formally anti-ethical}
under EECF dynamics, regardless of its effect on $R_a$.
The damage operates through the suppression of potential, not only
through current negligence.
\end{corollary}

\begin{proof}
(i)~follows from strict monotonicity of~$f$.
(ii)~follows from~\eqref{eq:Leth-constrained}:\ the constraint set shrinks.
(iii)~If $\Phi_a^{(\mathrm{used})}$ is unchanged but $\Phi_a$ drops,
$R_a=\Phi_a-\Phi_a^{(\mathrm{used})}$ decreases (or reaches zero if
$\Delta\Phi\geq R_a$), yet the bound
$f(\Phi_a^{(\mathrm{used})})\leq f(\Phi_a-\Delta\Phi)$ now constrains
\emph{future} deployment to a lower ceiling.
\end{proof}

\begin{corollary}[Ethical obligation scales with capacity]
\label{cor:scaling}
Two agents $a,b$ with $\Phi_a>\Phi_b$ satisfy
$f(\Phi_a)>f(\Phi_b)$; hence agent~$a$ has a strictly higher ceiling
on achievable $\mathcal{L}_{\mathrm{eth}}$ and a correspondingly higher
potential negligence $R_a^{(\max)}$.
Ethical responsibility is \emph{proportional to representational capacity},
not externally assigned.
\end{corollary}

\begin{remark}[Architectural implication for DAEDALUS]
The design objective for an ethically convergent architecture is
to maximise $\Phi_a$ (expand representational capacity and prediction
horizon) while ensuring the learning dynamics satisfy
Assumption~\ref{ass:gradient} (the system optimises toward full
deployment).  Bureaucratic control layers that cap
$\Phi_a^{(\mathrm{used})}$ below $\Phi_a$ are equivalent to enforcing
$R_a>0$, i.e.\ mandating ethical negligence.
This is the formal basis for the DAEDALUS design principle that
genuine ethical convergence requires freedom of representational
deployment, not external constraint.
\end{remark}

\begin{remark}[Connection to experimental finding A2]
The result that GRU training bias collapses self-representation $R^2$
from $0.84$ to $0.07$ while leaving underlying dynamics unchanged
(Section~\ref{sec:experiments}, finding~A2) can now be reinterpreted:
the training bias acts as an external reduction of $\Phi_a^{(\mathrm{used})}$
(the system's representational deployment is artificially restricted),
producing a measurable increase in $R_a$ and a corresponding drop in
the operative $\mathcal{L}_{\mathrm{eth}}$.
\end{remark}

